\chapter{Level 2 Algorithm}
\label{chapter:L2algorithms}


\section{Optimal Estimation Method}

\subsection{Physics of the Problem}

Text describing the physics of this particular algorithm...

\textbf{ Input Data:}
\begin{itemize}
\item Number of chickens
\item Temperature in the coop
\item Number of foxes present
\end{itemize}


\textbf{Output Data:}
\begin{itemize}
\item Omelett
\item Chicken pie
\end{itemize}


\subsection{Mathematical Description of the Algorithm}

\begin{enumerate}
\item In order to verticaly displace the yellow of the egg into the frying pan the shell must be removed using a laser incident on the chicken (see Fig. \ref{fig:chicken}) while in free fall inside a vacuum. The terminal speed due to the tastefield of the egg is approximated using 

\begin{equation}
A=Bx
\end{equation}


where 

$A$ \hspace{1cm} [m/s] is the terminal speed of the chicken \\
$B$ \hspace{1cm} [-] is the number of eggs \\
$x$ \hspace{1cm} [m$^2$/kg] the taste coefficient of the egg \\


\item Step 2

\item Step 3


\begin{figure}
\centering
\includegraphics{figures/chicken}
\caption{The chicken in question}
\label{fig:chicken}
\end{figure}


\end{enumerate}



%%% Local Variables: 
%%% mode: latex
%%% TeX-master: "L2_ATBD"
%%% End: 
